%document
\documentclass[10pt]{beamer}
%theme
\usetheme{metropolis}
% packages
\usepackage{color}
\usepackage{listings}
\usepackage[ngerman]{babel}
\usepackage[utf8]{inputenc}
\usepackage{multicol}

\usepackage{hyperref}
\hypersetup{colorlinks=true, linkcolor=blue, urlcolor=red}

% color definitions
\definecolor{mygreen}{rgb}{0,0.6,0}
\definecolor{mygray}{rgb}{0.5,0.5,0.5}
\definecolor{mymauve}{rgb}{0.58,0,0.82}

\lstset{
    backgroundcolor=\color{white},
    % choose the background color;
    % you must add \usepackage{color} or \usepackage{xcolor}
    basicstyle=\footnotesize\ttfamily,
    % the size of the fonts that are used for the code
    breakatwhitespace=false,
    % sets if automatic breaks should only happen at whitespace
    breaklines=true,                 % sets automatic line breaking
    captionpos=b,                    % sets the caption-position to bottom
    commentstyle=\color{mygreen},    % comment style
    % deletekeywords={...},
    % if you want to delete keywords from the given language
    extendedchars=true,
    % lets you use non-ASCII characters;
    % for 8-bits encodings only, does not work with UTF-8
    frame=single,                    % adds a frame around the code
    keepspaces=true,
    % keeps spaces in text,
    % useful for keeping indentation of code
    % (possibly needs columns=flexible)
    keywordstyle=\color{blue},       % keyword style
    % morekeywords={*,...},
    % if you want to add more keywords to the set
    numbers=left,
    % where to put the line-numbers; possible values are (none, left, right)
    numbersep=5pt,
    % how far the line-numbers are from the code
    numberstyle=\tiny\color{mygray},
    % the style that is used for the line-numbers
    rulecolor=\color{black},
    % if not set, the frame-color may be changed on line-breaks
    % within not-black text (e.g. comments (green here))
    stepnumber=1,
    % the step between two line-numbers.
    % If it's 1, each line will be numbered
    stringstyle=\color{mymauve},     % string literal style
    tabsize=4,                       % sets default tabsize to 4 spaces
    % show the filename of files included with \lstinputlisting;
    % also try caption instead of title
    language = C,
	showspaces = false,
	showtabs = false,
	showstringspaces = false,
	escapechar = ,
}

\def\ContinueLineNumber{\lstset{firstnumber=last}}
\def\StartLineAt#1{\lstset{firstnumber=#1}}
\let\numberLineAt\StartLineAt



\newcommand{\codeline}[1]{
	\alert{\texttt{#1}}
}

% This Document contains the information about this course.

% Authors of the slides
\author{
    Lecturers:
	%name1
	Mirko Jantschke,
	%name2
	Pascal Scholz
	\\
	%creaters of the slides
	Created by: Richard Mörbitz, Manuel Thieme
}


% Fancy Logo
\titlegraphic{\hfill\includegraphics[height=1.25cm]{../templates/fsr_logo_cropped}}


\title{C-Lessons}
\subtitle{Variables}
\date{\today}



% the actual document
\begin{document}

\maketitle

%-------------------------------------------------------------------------------------------------

\begin{frame}{Contents}
	\tableofcontents
\end{frame}

%-------------------------------------------------------------------------------------------------
\section{Bit operations}
%-------------------------------------------------------------------------------------------------

\begin{frame}{A little bit of logic}
	As you know, all data is stored as \textit{binary numbers} - sequences of 0 and 1.\\
	In C, you can operate on this bit layer by using the following \textit{logical bit} and \textit{shift} operators:\bigskip
	
	\begin{tabular}{|c|c|l|}
																						  	  \hline
		\textbf{Symbol} 	& \textbf{Operation} 	& \textbf{Example} 							\\\hline
		$|$					& logical or				& $0110\ |\ 0101 == 0111$ 		\\\hline
		$\&$ 				& logical and 				& $0110\ \&\ 0101 == 0100$ 	\\\hline
		\textasciicircum				& logical xor 				& $0110\ $\textasciicircum\ $0101 == 0011$ 	\\\hline
		\textasciitilde			& logical negation 			& \textasciitilde $0110 == 1001$	\\\hline
		$<<$ 			& shift to the left 			& $0110\ <<\ 2 == 011000$ 	\\\hline
		$>>$ 		& shift to the right 				& $0110\ >>\ 2 == 0001$ 			\\\hline
	\end{tabular}
	
\end{frame}

%-------------------------------------------------------------------------------------------------

\begin{frame}[fragile]{Computational arithmetics}
	With bit operations, some mathematical tasks can be solved more efficiently:
	\begin{itemize}
		\item Multiplying/dividing by $2^n$ is equivalent to a shift by $n$ bits
		\begin{lstlisting}
5 * 8 == 5 << 3;
60 / 4 == 60 >> 2;
\end{lstlisting}
		\item Instead of $\%\ 2^n$ you can use $\&\ 2^{n-1}$ 
		\begin{lstlisting}
22 % 2 == 22 & 1;
24 % 16 == 24 & 15;
\end{lstlisting}
	\end{itemize}\bigskip
	Be aware of the fact that the readability of your code will suffer from that. Most of these optimisations are done by the compiler anyway.
\end{frame}

%-------------------------------------------------------------------------------------------------

\begin{frame}[fragile]{Masking}
	$x\ |\ mask$ sets all bits in $x$ that are $1$ in $mask$.\\
	\begin{lstlisting}
'A' | 32;	/* Let a capital letter be small */
\end{lstlisting}\bigskip
	$x\ \&\ mask$ deletes all bits in $x$ that are $0$ in $mask$.\\
	\begin{lstlisting}
'a' & ~32;	/* Let a small letter be capital */
\end{lstlisting}\bigskip
	$x\ $\textasciicircum\ $mask$ inverts all bits in $x$ that are $1$ in $mask$.\\
	\begin{lstlisting}
'a' ^ 32;	/* "Toggle" a letter */
\end{lstlisting}
\end{frame}

%-------------------------------------------------------------------------------------------------

\begin{frame}[fragile]{Bit fields}
	Although it may seem efficient to use each bit of a number to store information in it, it will become nasty to access all the values by:\\
	$x\ \&\ 1$, $x\ \&\ 2$, $x\ \&\ 4$, \dots all the way up to $2^{sizeof\ int - 1}$\\
	\bigskip
	For this particular reason, C offers \textit{bit fields} like the following:
	\begin{lstlisting}
struct traffic_light {
	int red		: 1;
	int yellow	: 1;
	int green	: 1;
	int			: 5;	/* not in use */
};
\end{lstlisting}
	The members of bit fields can be accessed as if they were members of ordinary \textit{struct}s.

\end{frame}

%-------------------------------------------------------------------------------------------------

\end{document}
